\documentclass[11pt,a4paper]{ctexart}
\usepackage[margin=1in]{geometry}
\usepackage{booktabs}
\usepackage{longtable}
\usepackage{array}
\usepackage{hyperref}
\usepackage{enumitem}
\hypersetup{colorlinks=true,linkcolor=blue,urlcolor=blue}
\setlist[itemize]{nosep,leftmargin=1.5em}
\setlist[enumerate]{nosep,leftmargin=1.8em}

\title{Geant4-Agent User Guide and Local Deployment\\Geant4-Agent 使用方法与本地部署说明}
\author{}
\date{2026-03-07}

\begin{document}
\maketitle

\section*{English}

\subsection*{1. Scope}

This guide is for:

\begin{enumerate}
\item running the project locally with the web UI and multi-turn dialogue
\item reproducing the same setup on another Windows machine
\end{enumerate}

The current version is suitable for:

\begin{itemize}
\item local validation
\item internal testing
\item controlled user trials
\end{itemize}

It is not yet suitable for unrestricted public deployment.

\subsection*{2. Recommended Environment}

\begin{itemize}
\item OS: Windows 10 / 11
\item Python: 3.10+
\item GPU: optional for inference, recommended for local BERT training
\end{itemize}

\subsection*{3. Create a Virtual Environment}

\begin{verbatim}
python -m venv .venv
.\.venv\Scripts\Activate.ps1
python -m pip install --upgrade pip
\end{verbatim}

If PowerShell blocks script execution:

\begin{verbatim}
Set-ExecutionPolicy -Scope Process RemoteSigned
\end{verbatim}

\subsection*{4. Install Runtime Dependencies}

Minimum runtime install:

\begin{verbatim}
python -m pip install torch transformers safetensors
\end{verbatim}

Example CUDA 12.1 install:

\begin{verbatim}
python -m pip install torch --index-url https://download.pytorch.org/whl/cu121
python -m pip install transformers safetensors
\end{verbatim}

\subsection*{5. LLM Provider Configuration}

Configuration directory:

\begin{itemize}
\item \path{nlu/bert\_lab/configs/}
\end{itemize}

Example configs:

\begin{itemize}
\item \path{ollama\_config.json}
\item \path{openai\_compatible\_config.example.json}
\item \path{deepseek\_api\_config.example.json}
\end{itemize}

Recommended local private config:

\begin{itemize}
\item \path{nlu/bert\_lab/configs/deepseek\_api.local.json}
\end{itemize}

\subsection*{6. Local Model Directories}

Runtime model preflight checks:

\begin{itemize}
\item structure model directory
\item \texttt{ner} model directory
\end{itemize}

The current code searches for one structure model under:

\begin{itemize}
\item \path{nlu/bert\_lab/models/structure\_controlled\_v4c\_e1}
\item \path{nlu/bert\_lab/models/structure\_controlled\_v3\_e1}
\item \path{nlu/bert\_lab/models/structure\_controlled\_smoke}
\item \path{nlu/bert\_lab/models/structure\_opt\_v3}
\item \path{nlu/bert\_lab/models/structure\_opt\_v2}
\item \path{nlu/bert\_lab/models/structure}
\end{itemize}

The default NER directory is:

\begin{itemize}
\item \path{nlu/bert\_lab/models/ner}
\end{itemize}

Each model directory should at least contain:

\begin{itemize}
\item \texttt{config.json}
\item tokenizer assets such as \texttt{tokenizer.json}, \texttt{vocab.txt}, or equivalent
\end{itemize}

The repository does not ship trained weights by default.

\subsection*{7. Start the Local Web UI}

\begin{verbatim}
python ui\web\server.py
\end{verbatim}

Then open:

\begin{itemize}
\item \url{http://127.0.0.1:8088}
\end{itemize}

The startup log prints:

\begin{itemize}
\item current provider
\item current model
\item current base URL
\item model preflight status
\end{itemize}

\subsection*{8. Basic Local Usage Flow}

Recommended usage sequence:

\begin{enumerate}
\item start the web server
\item confirm the active LLM config
\item submit a natural-language request
\item answer follow-up questions until the configuration closes
\item inspect the user-facing summary and JSON output
\end{enumerate}

Currently supported interaction features:

\begin{itemize}
\item conservative handling of fuzzy first turns
\item overwrite confirm / reject
\item family-aware follow-up questions for graph-like geometry
\end{itemize}

\subsection*{9. Geometry-Only Theoretical Validation}

\begin{verbatim}
python -m builder.geometry.cli run_all --outdir builder/geometry/out --n_samples 200 --n_param_sets 100 --seed 7 --dataset builder/geometry/examples/coverage.csv
\end{verbatim}

Expected outputs:

\begin{itemize}
\item \path{coverage\_summary.json}
\item \path{coverage\_checked.jsonl}
\item \path{feasibility\_summary.json}
\item \path{ambiguity\_summary.json}
\end{itemize}

\subsection*{10. Regression Commands}

Stable validated live set:

\begin{verbatim}
python scripts/run_casebank_regression.py --dataset docs/eval_casebank_multiturn_live_v2_12.json --config nlu/bert_lab/configs/deepseek_api.local.json --outdir docs --baseline docs/casebank_baseline.json --min_confidence 0.6
\end{verbatim}

Expanded geometry-heavy live set:

\begin{verbatim}
python scripts/run_casebank_regression.py --dataset docs/eval_casebank_multiturn_live_v3_24.json --config nlu/bert_lab/configs/deepseek_api.local.json --outdir docs --baseline docs/casebank_baseline.json --min_confidence 0.6
\end{verbatim}

Dataset distribution check:

\begin{verbatim}
python scripts/verify_eval_casebank.py --dataset docs/eval_casebank_multiturn_live_v3_24.json
\end{verbatim}

\subsection*{11. Train the Structure Model Locally}

Build the normalized-text dataset:

\begin{verbatim}
python scripts/build_structure_v2_dataset.py --base nlu/bert_lab/data/controlled_structure.jsonl --extra nlu/bert_lab/data/controlled_multitask.jsonl --out nlu/bert_lab/data/controlled_structure_v2.jsonl --keep_original
\end{verbatim}

Train:

\begin{verbatim}
python scripts/train_structure_v2.py --config nlu/bert_lab/configs/structure_train_v2.json --outdir nlu/bert_lab/models/structure_controlled_v5_e2 --device cuda
\end{verbatim}

\subsection*{12. Deploy on Another Windows Machine}

Recommended sequence:

\begin{enumerate}
\item copy the repository
\item create \texttt{.venv}
\item install runtime dependencies
\item copy trained model directories into \path{nlu/bert\_lab/models/}
\item prepare one LLM config file
\item run \texttt{python ui\textbackslash web\textbackslash server.py}
\item open \url{http://127.0.0.1:8088}
\end{enumerate}

If the LLM runs remotely, only the config \path{base\_url} needs to point to the forwarded endpoint.

\subsection*{13. Deployment Recommendation}

Current recommendation:

\begin{itemize}
\item good for local use, internal testing, and controlled trials
\item not yet ready for unrestricted public deployment
\end{itemize}

\newpage
\section*{中文}

\subsection*{1. 适用范围}

这份说明面向两类场景：

\begin{enumerate}
\item 在本机启动 Web UI，运行多轮对话
\item 在另一台 Windows 机器上复现同样的环境
\end{enumerate}

当前版本适合：

\begin{itemize}
\item 本地验证
\item 内部测试
\item 受控真人试用
\end{itemize}

当前版本不适合直接做不受控公开部署。

\subsection*{2. 推荐环境}

\begin{itemize}
\item 操作系统：Windows 10 / 11
\item Python：3.10+
\item GPU：推理可选，训练 BERT 时建议使用 NVIDIA GPU
\end{itemize}

\subsection*{3. 创建虚拟环境}

\begin{verbatim}
python -m venv .venv
.\.venv\Scripts\Activate.ps1
python -m pip install --upgrade pip
\end{verbatim}

如果 PowerShell 限制脚本执行：

\begin{verbatim}
Set-ExecutionPolicy -Scope Process RemoteSigned
\end{verbatim}

\subsection*{4. 安装运行依赖}

最低运行安装：

\begin{verbatim}
python -m pip install torch transformers safetensors
\end{verbatim}

CUDA 12.1 示例：

\begin{verbatim}
python -m pip install torch --index-url https://download.pytorch.org/whl/cu121
python -m pip install transformers safetensors
\end{verbatim}

\subsection*{5. LLM 配置}

配置目录：

\begin{itemize}
\item \path{nlu/bert\_lab/configs/}
\end{itemize}

可用示例：

\begin{itemize}
\item \path{ollama\_config.json}
\item \path{openai\_compatible\_config.example.json}
\item \path{deepseek\_api\_config.example.json}
\end{itemize}

推荐本地私有配置：

\begin{itemize}
\item \path{nlu/bert\_lab/configs/deepseek\_api.local.json}
\end{itemize}

\subsection*{6. 本地模型目录}

运行时会检查：

\begin{itemize}
\item structure 模型目录
\item \texttt{ner} 模型目录
\end{itemize}

默认 structure 候选目录包括：

\begin{itemize}
\item \path{nlu/bert\_lab/models/structure\_controlled\_v4c\_e1}
\item \path{nlu/bert\_lab/models/structure\_controlled\_v3\_e1}
\item \path{nlu/bert\_lab/models/structure\_controlled\_smoke}
\item \path{nlu/bert\_lab/models/structure\_opt\_v3}
\item \path{nlu/bert\_lab/models/structure\_opt\_v2}
\item \path{nlu/bert\_lab/models/structure}
\end{itemize}

默认 NER 目录：

\begin{itemize}
\item \path{nlu/bert\_lab/models/ner}
\end{itemize}

模型目录至少应包含：

\begin{itemize}
\item \texttt{config.json}
\item tokenizer 资产，如 \texttt{tokenizer.json} 或 \texttt{vocab.txt}
\end{itemize}

仓库默认不附带训练好的权重。

\subsection*{7. 启动本地 Web UI}

\begin{verbatim}
python ui\web\server.py
\end{verbatim}

然后访问：

\begin{itemize}
\item \url{http://127.0.0.1:8088}
\end{itemize}

启动日志会显示：

\begin{itemize}
\item 当前 provider
\item 当前 model
\item 当前 base URL
\item model preflight 状态
\end{itemize}

\subsection*{8. 本地使用流程}

推荐顺序：

\begin{enumerate}
\item 启动 Web 服务
\item 确认当前 LLM 配置
\item 输入自然语言需求
\item 按系统追问补齐信息
\item 查看用户层摘要和 JSON 输出
\end{enumerate}

当前已经支持：

\begin{itemize}
\item fuzzy 首轮的保守处理
\item overwrite 确认 / 拒绝
\item graph 几何的 family-aware 追问
\end{itemize}

\subsection*{9. 仅运行 geometry 理论验证}

\begin{verbatim}
python -m builder.geometry.cli run_all --outdir builder/geometry/out --n_samples 200 --n_param_sets 100 --seed 7 --dataset builder/geometry/examples/coverage.csv
\end{verbatim}

主要输出：

\begin{itemize}
\item \path{coverage\_summary.json}
\item \path{coverage\_checked.jsonl}
\item \path{feasibility\_summary.json}
\item \path{ambiguity\_summary.json}
\end{itemize}

\subsection*{10. 回归命令}

当前稳定通过的 live 集：

\begin{verbatim}
python scripts/run_casebank_regression.py --dataset docs/eval_casebank_multiturn_live_v2_12.json --config nlu/bert_lab/configs/deepseek_api.local.json --outdir docs --baseline docs/casebank_baseline.json --min_confidence 0.6
\end{verbatim}

扩展的 geometry-heavy live 集：

\begin{verbatim}
python scripts/run_casebank_regression.py --dataset docs/eval_casebank_multiturn_live_v3_24.json --config nlu/bert_lab/configs/deepseek_api.local.json --outdir docs --baseline docs/casebank_baseline.json --min_confidence 0.6
\end{verbatim}

检查数据集分布：

\begin{verbatim}
python scripts/verify_eval_casebank.py --dataset docs/eval_casebank_multiturn_live_v3_24.json
\end{verbatim}

\subsection*{11. 本地训练 structure 模型}

构建数据集：

\begin{verbatim}
python scripts/build_structure_v2_dataset.py --base nlu/bert_lab/data/controlled_structure.jsonl --extra nlu/bert_lab/data/controlled_multitask.jsonl --out nlu/bert_lab/data/controlled_structure_v2.jsonl --keep_original
\end{verbatim}

训练命令：

\begin{verbatim}
python scripts/train_structure_v2.py --config nlu/bert_lab/configs/structure_train_v2.json --outdir nlu/bert_lab/models/structure_controlled_v5_e2 --device cuda
\end{verbatim}

\subsection*{12. 在另一台 Windows 机器部署}

推荐顺序：

\begin{enumerate}
\item 复制仓库
\item 创建 \texttt{.venv}
\item 安装运行依赖
\item 复制训练好的模型到 \path{nlu/bert\_lab/models/}
\item 准备一个 LLM 配置文件
\item 运行 \texttt{python ui\textbackslash web\textbackslash server.py}
\item 打开 \url{http://127.0.0.1:8088}
\end{enumerate}

如果 LLM 在远端，只需把配置中的 \path{base\_url} 指向 SSH 转发后的地址。

\subsection*{13. 部署建议}

当前建议是：

\begin{itemize}
\item 适合本地使用、内部测试、受控试用
\item 不建议直接进行不受控公开部署
\end{itemize}

\end{document}
